
%%%%%%%%%%%%%%%%%%%%%%%%%%%%%%%%%%%%%%%%%%%%%%%%%%%
\section{UV Plane}
% ------------------------------------------------
\SC{Add cartoon}



The projected baseline between two antennas is decomposed into two coordinates $ (u, v) $ with units of observing wavelength along the east and north directions (respectively) as projected in the direction of the target \citep{Monnier}. For ground based observing (like ALMA), the object moves across the sky along a line of constant declination, so the $(u, v)$ coverage changes with time \citep{Monnier}. This is actually a benefit since Earth's rotation dramatically increases the $(u,v)$ coverage without requiring additional telescopes \citep{Monnier}. 

% The $u,v$ coverage depends on the projected baseline separations in the direction of the target (not the physical baselines) \citep{Monnier}. 



Each pair of telescopes measures one complex visibility, which describes the amplitude and phase of the interference pattern, as a function of time and frequency \citep{Ricci2017}.  Each complex visibility is a single Fourier component of the sky brightness distribution, and together they sample its Fourier transform \citep{Marvil2026, Monnier}. The visibility function  $V(u, v) $ is given by:
\begin{equation}
    V(u, v) = \int_{-\infty}^{\infty} I(l, m) e^{-2\pi i(ul + vm)} dl dm
    \label{eq: V(u,v)}
\end{equation}
\noindent where $I(l, m)$ is the intensity function, and $(l, m)$ are the angular sky coordinates  \citep{Marvil2026}. 












%%%%%%%%%%%%%%%%%%%%%%%%%%%%%%%%%%%%%%%%%%%%%%%%%%%
% \quoted{Here $l$ and $m$ have units of radians} \citep{Monnier}. 

% A major difference bweteen a single-telescope observation and observations from an interderometer is the (u,v) coverage \citep{Monnier}.


% A visibility is a measurement of one Fourier component of the sky brightness distribution, made in the uv space \citep{Marvil2026}. 

% \quoted{The vector $\vec{sigma} = (l,m)$ can be represented in rectilinear coorinates on the celestial sphere, where $l$ points along local east and $m$ points north$^2$}\citep{Monnier}.

% The $(u, v)$ axes lie in a plane paralle to the Earth's equator \citep{Thompson}. 



% Gridding combines all the measured visibilities onto a UV grid \citep{Marvil2026}. 

% A baseline and its conjugate trace out elliptical arcs through the uv plane over time beacuse of earths rotation \citep{Marvil2026}. 


 


% The easy-west antenna spacings contain components in this plane only, and as the Earth rotates, the specing vectors sweep out circles with the $(u,v)$ origin \citep{Thompson}. 

%%%%%%%%%%%%%%%%%%%%%%%%%%%%%%%%%%%%%%%%%%%%%%%%%%%

